\subsection{Summary and Contributions}

This thesis developed a workflow for studying 5G mobility from two complementary viewpoints: radio-level measurement behavior and protocol-level handover execution. The Python radio simulator generates structured mobility and radio-measurement traces, including UE position, UE height, serving-cell information, RSRP, SINR, shadow fading, and handover indicators. The StormSim-based emulation workflow then replays measurement-driven serving-cell changes in a free5GC environment and records handover trial summaries and timer events.

The main contribution of the thesis is the separation between a radio-level serving-cell change and a completed protocol-level handover. In the radio simulator, a handover indicator means that the selected serving gNB changes according to an RSRP-based rule. In StormSim, the handover is treated as a protocol procedure involving source gNB, target gNB, UE context transfer, path switch, T304, and Xn relocation timers. This separation avoids treating a measurement decision as if it were already a completed handover.

A second contribution is the construction of a structured radio-data generation pipeline for terrestrial and aerial UE-height ranges. The simulator supports UMa and UMa-AV propagation behavior, height-dependent mobility settings, hexagonal gNB placement, and CSV output suitable for post-processing. A third contribution is the creation of a StormSim-ready seven-gNB measurement trace with fixed gNB-indexed RSRP columns, allowing the generated radio data to be replayed by the handover emulator without losing the identity of each gNB.

\subsection{Main Findings}

The radio-measurement results show that UE height has different effects on RSRP and SINR. Compared with the terrestrial height ranges, the aerial height ranges produce stronger average serving-cell RSRP because the UE has more favorable propagation and a higher probability of LOS connection. However, the SINR does not improve in the same way. At aerial heights, the UE can also receive stronger signals from neighbouring gNBs, so co-channel interference increases. As a result, the average SINR decreases in the H2 and H3 altitude ranges even though the serving-cell RSRP is higher. This confirms that RSRP alone is not sufficient for evaluating A2G mobility quality.

The generated radio traces also show that shadow fading introduces slow signal variation around the path-loss trend. The shadow-fading distribution becomes narrower at higher altitude, especially in the H3 range. This behavior is consistent with the implemented height-dependent propagation model and helps explain why aerial traces can have stronger but more interference-limited received signals.

For protocol-level evaluation, a seven-gNB measurement trace was generated with one center gNB and six surrounding first-tier gNBs. The trace contains 301 measurement steps and 18 serving-cell changes. When replayed in StormSim as Xn handover opportunities under the baseline no-loss condition, all 18 handover trials completed successfully. The successful trial-summary delay has a mean of 200.530 ms, a median of 200.550 ms, and a 95th percentile of 201.013 ms.

The timer-event log provides a finer interpretation of this delay. In the same run, \texttt{T304} stopped successfully for all 18 handovers with an average elapsed time of 10.621 ms, while \texttt{TXnRELOCprep} and \texttt{TXnRELOCoverall} stopped within about 1 ms on average. Therefore, the approximately 200 ms trial-summary delay mainly reflects the scenario-level polling interval used to observe handover completion, not the duration of every internal protocol phase.

\subsection{Limitations}

The radio simulator is a controlled system-level model rather than a full physical-layer implementation. RSRP is derived from received power using a resource-grid approximation, and the serving-cell decision uses an RSRP margin without time-to-trigger, measurement filtering, or full RRC measurement-event processing. The LOS/NLOS state and shadow fading are cached per UE--gNB link to keep the trace spatially stable, but this remains an abstraction of real propagation dynamics.

The current StormSim result is a baseline validation of the measurement-driven workflow, not a complete performance study under all network conditions. The verified seven-gNB run uses Xn handover with zero packet loss, zero added latency, and zero jitter. It demonstrates that the generated measurement trace can drive successful protocol-level handovers, but it does not yet quantify handover robustness under degraded transport conditions.

The trial-summary delay is also affected by the implementation of the scenario runner. Completion is checked using a 200 ms polling ticker, so the exported successful handover duration is quantized around 200 ms. Timer-event logs are therefore necessary when interpreting lower-level handover phases. In addition, the CSV interface between the Python simulator and StormSim still requires a dedicated adapter script for fixed gNB-indexed RSRP columns.

Finally, the present evaluation focuses on Xn handover. N2 handover, full CHO mode evaluation, RLF behavior, beam-failure behavior, and systematic loss--delay experiments remain outside the verified result set reported in this thesis.

\subsection{Future Work}

Future work should first unify the radio-simulator and StormSim CSV formats so that fixed gNB-indexed RSRP fields are produced directly by the main simulator logger. This would remove the need for a separate adapter script and make measurement-driven handover experiments easier to reproduce.

The handover-delay evaluation should then be extended from the baseline no-loss case to a controlled matrix of packet loss, latency, and jitter settings. This matrix should report success rate, failure reason, timeout behavior, mean delay, median delay, and percentile delay for each setting. Xn and N2 handover should also be compared under the same measurement trace so that the effect of direct gNB-to-gNB handover and core-network-assisted handover can be separated.

Another important direction is to reduce or replace the 200 ms polling interval in the scenario runner. Event-driven completion logging would allow the trial-summary delay to better match the internal handover completion time. The timer lifecycle should also be checked so that each handover trial has complete and consistent start and terminal timer records.

On the radio side, future work can add time-to-trigger, measurement filtering, and additional RRC event conditions such as A5 to make the serving-cell and handover-trigger logic closer to real mobility procedures. The generated dataset can also be expanded to multiple UE trajectories and repeated fixed-seed scenarios to support more systematic statistical comparison.

\subsection{Chapter Conclusion}

This chapter summarized the thesis contributions, main findings, limitations, and future research directions. Overall, the thesis provides a practical platform for generating 5G mobility traces and analyzing handover behavior beyond simple serving-cell changes. The current results show that aerial UE height can improve RSRP while reducing SINR due to interference, and that a generated seven-gNB trace can drive successful Xn handovers in StormSim under baseline no-loss conditions. Broader claims about handover robustness require the controlled network-impairment and handover-type experiments described as future work.
